\section{Lean verification of Theorem A}\label{app:leanA}
Theorem A is proved in Lean~4, conditional only on
\nolinkurl{UpperRealArithmetic.TarskiDecision}: truth of existentially
quantified rational ordered-ring sentences in disjunctive normal form is
computable. This is a special case of Tarski's decision theorem for real
arithmetic \cite[Section~2, Theorem~31 and Corollary~32, p.~39]{Tarski1951}.
It is an explicit argument of the theorem; the decision algorithm is not
formalized. The comparator challenge \nolinkurl{lean/Challenge.lean}
states Theorem A as follows, with \nolinkurl{TheoremAStatement} as defined
in \nolinkurl{lean/Sigma2/Statements.lean}:
\begin{flushleft}\leftskip1.5em\footnotesize\ttfamily\frenchspacing
\detokenize{theorem theoremA_of_tarski_paper (hT : UpperRealArithmetic.TarskiDecision) :}\\
\hspace*{2em}\detokenize{TheoremAStatement}\\[4pt]
\detokenize{def TheoremAStatement : Prop :=}\\
\hspace*{1em}\detokenize{SigmaTwoComplete SpatialPerfect} $\wedge$ \detokenize{SigmaTwoComplete EntropyPerfect} $\wedge$\\
\hspace*{1em}\detokenize{PromiseHard SpatialPerfect PromisedGame} $\wedge$ \detokenize{PromiseHard EntropyPerfect PromisedGame}
\end{flushleft}
Here the languages consist of natural-number codes of finite rational
Boolean two-player games, and invalid codes belong to none of them.
\texttt{SpatialPerfect} is the set of games with a perfect normal spatial
strategy (perfection is attained, not merely value one), and
\texttt{EntropyPerfect} the set of games with such a strategy whose two
reduced density operators have finite entropy; Section~\ref{sec:prelim}
explains why the one-marginal entropy convention gives the same language.
\texttt{SigmaTwoComplete} means $\Sigma^0_2$ membership and hardness under
total computable many-one reductions. \texttt{PromiseHard} says that each
hardness reduction can be chosen so that every output lies in
\texttt{PromisedGame}: it has finite-dimensional value one and no perfect
finite-dimensional strategy. The same statement is proved as
\nolinkurl{theoremA_of_tarski} in \nolinkurl{lean/Sigma2/Main.lean} and as
\nolinkurl{theoremA_of_tarski_paper}, through the printed admissibility
predicate, in \nolinkurl{lean/Sigma2/MainPaper.lean}. The theorem
\nolinkurl{theoremA_shared_family} in \nolinkurl{lean/Sigma2/SharedFamily.lean}
proves that a single total computable map is a promised reduction for both
languages, which is the family of Theorem A.

\subsection{What is checked}
The first column uses this paper's numbering. The Lean names
\nolinkurl{lemma_4_4_strong_of_obligations} and
\nolinkurl{lemma_3_1_compact_sandwiches} refer to Lemma~\ref{lem:four} at
$n=4$ and to Lemma~\ref{lem:sandwich}. Several theorems in the table take
as an argument a compiler specification, of type
\nolinkurl{CompilerSpecification.Specification}. The final theorems supply
it from the proved obligations \nolinkurl{CompilerObligations.obligations_all},
so it is not a hypothesis of any final theorem. The only hypothesis of the final theorems is
\nolinkurl{TarskiDecision}; the rows below say where it enters.

\begin{center}\footnotesize
\begin{minipage}{\linewidth}
\setlength{\tabcolsep}{4pt}
\begin{tabular}{@{}>{\raggedright\arraybackslash}p{.17\linewidth}>{\raggedright\arraybackslash}p{.44\linewidth}>{\raggedright\arraybackslash}p{.33\linewidth}@{}}
\toprule Printed result & Lean name & Checked form\\\midrule
Theorem A &
\nolinkurl{Sigma2.theoremA_of_tarski};
\nolinkurl{Sigma2.theoremA_of_tarski_paper};
\nolinkurl{Sigma2.theoremA_shared_family} &
Full, given \texttt{TarskiDecision} only: both completeness clauses, both
promised-hardness clauses, and one computable hardness family for both
languages (the third name, which needs no hypothesis).\\[4pt]
Lemma 4.5, $n=4$ &
\nolinkurl{Sigma2.CompilerContract.lemma_4_4_strong_of_obligations};
\nolinkurl{Sigma2.CompilerSpecification.strong_compiler};
\nolinkurl{Sigma2.CompilerSpecification.four_bridge_compiler};
\nolinkurl{Sigma2.CompilerObligations.obligations_all} &
Partial: four-bridge exact compilation, with all seven compiler
obligations proved; not general $n$ or commuting hosts.\\[4pt]
Theorem 3.1, membership clauses; Lemma 3.2 &
\nolinkurl{Sigma2.Upper.uniform_strata_piOne_of_tarski};
\nolinkurl{Sigma2.Upper.spatial_union};
\nolinkurl{Sigma2.Upper.spatial_sigmaTwo_of_tarski};
\nolinkurl{Sigma2.Upper.lemma_3_1_compact_sandwiches} &
Partial: the first two clauses of Theorem 3.1, with the strata and
$\Sigma^0_2$ membership given \texttt{TarskiDecision}; Lemma 3.2.\\[4pt]
Proposition 3.3; Corollary 3.4; equation (5) &
\nolinkurl{Sigma2.UpperEntropyAssembly.perfect_of_sequence};
\nolinkurl{Sigma2.Upper.entropy_tail};
\nolinkurl{Sigma2.Upper.entropy_stage_slack};
\nolinkurl{Sigma2.Upper.entropy_scheduled_iff};
\nolinkurl{Sigma2.Upper.entropy_relaxation_computable_of_tarski};
\nolinkurl{Sigma2.Upper.entropy_sigmaTwo_of_tarski} &
Partial: perfect-game compactness and effective finite-entropy
membership, the latter given \texttt{TarskiDecision}; not
prescribed-table compactness, entropy infima or optimal entropy.\\[4pt]
Sections 2 and 4, normal toolkit &
\nolinkurl{Sigma2.Preliminaries.normal_pure_perfect_equivalence} &
Partial: the normal/pure perfection equivalence, not all of
(N1)--(N8).\\[4pt]
Lemma 4.1 &
\nolinkurl{Sigma2.Compiler.lemma_4_1_normal_bridge};
\nolinkurl{Sigma2.Compiler.calibrated_bridge_vector} &
Partial: the spatial support interface and component identity; not
commuting hosts.\\
\bottomrule
\end{tabular}
\end{minipage}
\end{center}

\begin{center}\footnotesize
\begin{minipage}{\linewidth}
\setlength{\tabcolsep}{4pt}
\begin{tabular}{@{}>{\raggedright\arraybackslash}p{.17\linewidth}>{\raggedright\arraybackslash}p{.44\linewidth}>{\raggedright\arraybackslash}p{.33\linewidth}@{}}
\toprule Printed result & Lean name & Checked form\\\midrule
Lemma 4.2 &
\nolinkurl{Sigma2.Compiler.lemma_4_2_label_refinement};
\nolinkurl{Sigma2.Compiler.joint_refinement_iff} &
Coarse-sign identification and the converse commuting refinement.\\[4pt]
Lemma 4.3; equation (19) &
\nolinkurl{Sigma2.Compiler.lemma_4_3_boolean_rows};
\nolinkurl{Sigma2.Compiler.row_loss_formula} &
Conditional: perfect games, with the \texttt{RowsHold} and
\texttt{NamedBinaryProjections} premises; (19) only as a trace
identity under a slot-sum hypothesis, not an event probability.\\[4pt]
Definition 4.4 &
\nolinkurl{Sigma2.Presentation};
\nolinkurl{Sigma2.Presentation.admissible};
\nolinkurl{Sigma2.NormalRealization};
\nolinkurl{Sigma2.Compiler.Compiled};
\nolinkurl{Sigma2.Compiler.ApproxCompiled};
\nolinkurl{Sigma2.PaperAdmissible};
\nolinkurl{Sigma2.PaperAdmissibility.application_paper_admissible};
\nolinkurl{Sigma2.PaperAdmissibility.saturation_bridge} &
\texttt{Presentation.admissible}, the domain of the compiler theorems,
omits two transport conditions of the printed definition.
\texttt{PaperAdmissible} is the printed definition; it holds for the
Theorem A input tables, and every admissible presentation saturates to a
\texttt{PaperAdmissible} one with equivalent normal realizations and the
same compiled game.\\[4pt]
Section 4.7, equations (28)--(29) &
\nolinkurl{Sigma2.CompilerSpecification.approximate_compilation};
\nolinkurl{Sigma2.CompilerApproxBase.approxCompiled_of_base15} &
Partial: an existential dimension-independent loss bound under exact
structural hypotheses; not the explicit coefficient (29), robust
extraction, or preservation of an input witness or its finite
dimension.\\[4pt]
Section 5.1, equations (30)--(33) &
\nolinkurl{Sigma2.Tree.fock_decomposition};
\nolinkurl{Sigma2.Tree.fock_density};
\nolinkurl{Sigma2.Tree.fock_entropy};
\nolinkurl{Sigma2.Tree.no_finite_realization} &
Rooted Fock structure, weights, entropy and finite-dimensional
obstruction.\\[4pt]
Lemma 5.1 &
\nolinkurl{Sigma2.Tree.lemma_5_1_definite_computation};
\nolinkurl{Sigma2.ReductionSoundBridge.nonreject_of_data} &
Classical tableau on a root-clock sector, with its concrete data from
the row-extraction modules.\\
\bottomrule
\end{tabular}
\end{minipage}
\end{center}

\begin{center}\footnotesize
\begin{minipage}{\linewidth}
\setlength{\tabcolsep}{4pt}
\begin{tabular}{@{}>{\raggedright\arraybackslash}p{.17\linewidth}>{\raggedright\arraybackslash}p{.44\linewidth}>{\raggedright\arraybackslash}p{.33\linewidth}@{}}
\toprule Printed result & Lean name & Checked form\\\midrule
Section 6, equations (40)--(46); Lemma 6.1 &
\nolinkurl{Sigma2.Clock.finite_phase_sector};
\nolinkurl{Sigma2.Clock.coordinate_clock};
\nolinkurl{Sigma2.Clock.exact_return};
\nolinkurl{Sigma2.Clock.lemma_6_1_good_primitive_phase};
\nolinkurl{Sigma2.Clock.bounded_reader} &
Conditional: \nolinkurl{finite_phase_sector}, \nolinkurl{coordinate_clock} and \nolinkurl{exact_return} take \nolinkurl{hP : P.admissible} and \nolinkurl{D : Clock.Data R}; the last two also take a root vector, and exact return a primitive-phase eigenvector. \nolinkurl{bounded_reader} assumes a primitive phase and its separation bound; Lemma~6.1 supplies a suitable phase and bound for every prime. \nolinkurl{Reduction.realization_sound} supplies the clock data for Theorem~A.\\[4pt]
Section 7; Theorem 7.1 &
\nolinkurl{Sigma2.Reduction.effective_program};
\nolinkurl{Sigma2.Reduction.application_admissible};
\nolinkurl{Sigma2.Reduction.presentation_codes_computable};
\nolinkurl{Sigma2.Reduction.realization_sound};
\nolinkurl{Sigma2.Reduction.realization_complete};
\nolinkurl{Sigma2.CompiledReduction.theorem_7_1_reduction} &
Effective input construction, both perfection equivalences and absence
of any finite-dimensional perfect output.\\[4pt]
Section 7.1, equations (50)--(54) &
\nolinkurl{Sigma2.CompiledFockApproximation.normal_approximation};
\nolinkurl{Sigma2.CompiledFockApproximation.finite_approximation};
\nolinkurl{Sigma2.CompiledValue.four_defect_approximation};
\nolinkurl{Sigma2.CompiledValue.value_one};
\nolinkurl{Sigma2.CompiledValue.reduction_promised} &
Partial: geometric approximation and the promise on every input; the
$Q,L,d$ of \nolinkurl{four_defect_approximation} are existential, not the
actual row counts or denominators.\\[4pt]
Other local results &
Thin core, \nolinkurl{lean/THIN-CORE.md} &
Local statements only; see Table~\ref{tab:leancore}.\\
\bottomrule
\end{tabular}
\end{minipage}
\end{center}

\noindent\textbf{Not formalized:} Theorems B--D beyond their thin-core
lemmas, including the general semifinite and type classification and the
communication classification; Theorem~\ref{thm:bell} beyond its
winning-functional identities, and the last two clauses of
Theorem~\ref{thm:strata}; the general $n$-bridge and commuting-host
compiler; the economical ten-label schedule, its size counts and
equation~\eqref{eq:valueconstant} with actual row denominators (the
formal compiler transports all typed labels); the reduction of strategies
on arbitrary Hilbert spaces to $\ell^2(\mathbb N)$ by compression;
Tarski's decision algorithm.

\noindent\textbf{Conventions:} both local spaces are $\ell^2(\mathbb N)$,
and normal states are countable ensembles of unit vectors in
$\ell^2(\mathbb N\times\mathbb N)$, their completed tensor product. A
finite-dimensional strategy is a unit vector supported on finitely many
basis vectors, with POVMs on $\ell^2(\mathbb N)$. Every game has at least
one question and one answer for each party.

\subsection{Trust base}
The proofs use Lean 4.35.0 and the Mathlib commit
\begin{quote}\small
\nolinkurl{56a005855ea17e0b0958f73ec801ae140b0a1679};
\end{quote}
every pin is in \nolinkurl{lean/}. Every declaration of the
\texttt{Sigma2} library, including every theorem in the table, uses no
axioms beyond Lean's three standard axioms:
\begingroup\scriptsize
\begin{verbatim}
'Sigma2.theoremA_of_tarski' depends on axioms: [propext, Classical.choice, Quot.sound]
'Sigma2.theoremA_of_tarski_paper' depends on axioms: [propext, Classical.choice, Quot.sound]
'Sigma2.theoremA_shared_family' depends on axioms: [propext, Classical.choice, Quot.sound]
\end{verbatim}\endgroup
The proof of Theorem A passes two independent kernel checkers,
\texttt{leanchecker} and \texttt{nanoda}. The comparator challenge
\nolinkurl{lean/Challenge.lean} imports only Mathlib and defines the
statement vocabulary transparently. It states Theorem A as displayed
above, the shared hardness family \nolinkurl{theoremA_shared_family}, and
the printed admissibility definition \nolinkurl{paper_admissible_iff}.
\nolinkurl{lean/Solution.lean} imports the proofs, and
\nolinkurl{lean/comparator.json} permits only the three standard axioms.
From \texttt{lean/}, on Linux with bubblewrap and with the dependencies
placed as below, one command checks all three:
\texttt{lake comparator -{}-config comparator.json}.
Records of the comparator run and the kernel checks are in
\nolinkurl{lean/evidence/}.

\subsection{Reproduce}
With the Lean 4.35.0 toolchain first on the path and the pinned
dependency sources placed under \nolinkurl{lean/vendor/} as described in
\nolinkurl{lean/README.md}, run from \texttt{lean/} (the first line is
PowerShell syntax):
\begin{flushleft}\leftskip1.5em\ttfamily\frenchspacing
\$env:LEAN\_NUM\_THREADS = \textquotesingle 1\textquotesingle\\
lake exe cache get\\
lake build\\
lake env lean -{}j1 -{}M6144 Audit.lean
\end{flushleft}
The first line selects one Lean thread, and \texttt{lake exe cache get}
fetches the Mathlib build cache. \texttt{lake build} rebuilds every project
module from source. It audits the axioms of every declaration of the
library, including private declarations, and rejects admitted proofs and
additional axioms; it audits the proofs of the two final Theorem A theorems and
checks that every Lean name in the table exists. The last command prints
the axioms of every theorem in the table; \nolinkurl{lean/AXIOMS.txt} is
the expected output.
